\documentclass[10pt]{extarticle}
\usepackage[a4paper,margin=1.5cm]{geometry}
\usepackage[T1]{fontenc}
\usepackage{amsmath,amssymb}
\usepackage{graphicx}
\usepackage[export]{adjustbox}
\usepackage{multicol}
\usepackage{enumitem}
\usepackage{booktabs,array,tabularx}
\setlength{\columnsep}{0.7cm}
\setlength{\parindent}{0pt}
\setlength{\parskip}{3pt}
\pagestyle{empty}
\makeatletter
% figure registry (figdefs.tex is generated by the extract step)
\newcommand{\qfigdef}[4]{\@namedef{qf@file@#1}{#2}\@namedef{qf@w@#1}{#3}\@namedef{qf@wide@#1}{#4}}
\newcommand{\qf@check}[1]{\@ifundefined{qf@file@#1}{\PackageError{score2}{Unknown figure id #1}{}}{}}
\newcommand{\qf@img}[4][]{\includegraphics[width=#2\dimexpr\@nameuse{qf@w@#4}pt\relax,max width=#3,max height=0.45\textheight,#1]{\@nameuse{qf@file@#4}}}
% Full-width questions: finalize wraps every question that contains a wide figure in \qfullbegin...\qfullend,
% so the WHOLE question (stem, figures, options) is set across the page in reading order, then the two columns
% resume. The question's head - its number, stem and first figure (\qfullsplit marks the end of the head) -
% is kept on one page, so a number is never left at a page foot with its figure overleaf; the rest may break
% across pages normally, so no big gaps are left. Without a \qfullsplit, the whole question is kept together
% if it is under 45% of a page. Never write these yourself.
\newif\ifqf@full
\newif\ifqf@boxed
\newsavebox{\qf@box}
\newcommand{\qf@open}{\global\qf@fulltrue\par\global\qf@boxedtrue
  \setbox\qf@box\vbox\bgroup\hsize\textwidth\linewidth\textwidth}
\newcommand{\qf@emit}[1]{\par\egroup\global\qf@boxedfalse
  \ifdim\dimexpr\ht\qf@box+\dp\qf@box\relax<#1\textheight\noindent\usebox{\qf@box}\par\else\unvbox\qf@box\fi}
\newcommand{\qfullsplit}{\ifqf@boxed\qf@emit{0.85}\fi}
\newcommand{\qf@close}{\ifqf@boxed\qf@emit{0.45}\fi\global\qf@fullfalse}
\newcommand{\qfullbegin}{\end{multicols}\qf@open}
\newcommand{\qfullend}{\qf@close\begin{multicols}{2}}
\newcommand{\qfullstart}{\qf@open}
\newcommand{\qfulldone}{\qf@close}
% \qfig[scale]{id}: block figure, centred, directly where it is written (size set by the tool).
% Fallback if a wide figure is not inside a full-width question: break out of the columns in place - no floats.
\newcommand{\qfig}[2][1]{\qf@check{#2}%
  \ifqf@full\qf@blk{#1}{#2}\else
  \ifnum\@nameuse{qf@wide@#2}=1 \expandafter\@firstoftwo\else\expandafter\@secondoftwo\fi
  {\end{multicols}\begin{center}\qf@img{#1}{\textwidth}{#2}\end{center}\begin{multicols}{2}}%
  {\qf@blk{#1}{#2}}\fi}
\newcommand{\qf@blk}[2]{\par\smallskip{\centering\qf@img{#1}{\linewidth}{#2}\par}\smallskip}
% \qfigi[scale]{id}: inline figure (inside option lists / tables / minipages); never breaks columns.
\newcommand{\qfigi}[2][1]{\qf@check{#2}\qf@img[valign=c]{#1}{\linewidth}{#2}}
\makeatother
% structure macros - lint reads these to check completeness and figure placement
\newcommand{\papersec}[1]{\par\bigskip{\centering\Large\bfseries PAPER-#1\par}\medskip}
\newcommand{\partsec}[2]{\par\medskip\noindent{\large\bfseries PART-#1: #2}\par\smallskip}
\newcommand{\sect}[1]{\par\smallskip\noindent\textbf{#1}\par}
\newcommand{\qn}[1]{\par\smallskip\noindent\textbf{#1.}~}
\newcommand{\paraq}[2]{\par\smallskip\noindent\textbf{Paragraph for Questions #1 and #2}\par}
\newenvironment{qoptions}{\begin{itemize}[leftmargin=1.8em,itemsep=1pt,topsep=2pt,label={}]}{\end{itemize}}
\newcommand{\opt}[2]{\item[(#1)] #2}
% figure registry (generated by score2_tools.py finalize)
\qfigdef{p1_phy_s2_q06_1}{Score Advanced - Test 38 (13 Apr 2025)-img/p1_phy_s2_q06_1.png}{141.0}{0}
\qfigdef{p1_phy_s2_q07_1}{Score Advanced - Test 38 (13 Apr 2025)-img/p1_phy_s2_q07_1.png}{125.3}{0}
\qfigdef{p1_phy_s3_q08_1}{Score Advanced - Test 38 (13 Apr 2025)-img/p1_phy_s3_q08_1.png}{127.8}{0}
\qfigdef{p1_phy_s3_q09_1}{Score Advanced - Test 38 (13 Apr 2025)-img/p1_phy_s3_q09_1.png}{119.3}{0}
\qfigdef{p1_phy_s3_q11_1}{Score Advanced - Test 38 (13 Apr 2025)-img/p1_phy_s3_q11_1.png}{151.9}{0}
\qfigdef{p1_phy_s4_q02_1}{Score Advanced - Test 38 (13 Apr 2025)-img/p1_phy_s4_q02_1.png}{163.4}{0}
\qfigdef{p1_phy_s4_q03_1}{Score Advanced - Test 38 (13 Apr 2025)-img/p1_phy_s4_q03_1.png}{125.4}{0}
\qfigdef{p1_phy_s4_q04_1}{Score Advanced - Test 38 (13 Apr 2025)-img/p1_phy_s4_q04_1.png}{156.5}{0}
\qfigdef{p1_phy_s4_q06_1}{Score Advanced - Test 38 (13 Apr 2025)-img/p1_phy_s4_q06_1.png}{153.4}{0}
\qfigdef{p1_chem_s1_q02_1}{Score Advanced - Test 38 (13 Apr 2025)-img/p1_chem_s1_q02_1.png}{216.5}{0}
\qfigdef{p1_chem_s1_q02_A}{Score Advanced - Test 38 (13 Apr 2025)-img/p1_chem_s1_q02_A.png}{105.4}{0}
\qfigdef{p1_chem_s1_q02_B}{Score Advanced - Test 38 (13 Apr 2025)-img/p1_chem_s1_q02_B.png}{106.4}{0}
\qfigdef{p1_chem_s1_q02_C}{Score Advanced - Test 38 (13 Apr 2025)-img/p1_chem_s1_q02_C.png}{105.9}{0}
\qfigdef{p1_chem_s1_q02_D}{Score Advanced - Test 38 (13 Apr 2025)-img/p1_chem_s1_q02_D.png}{106.3}{0}
\qfigdef{p1_chem_s2_q06_1}{Score Advanced - Test 38 (13 Apr 2025)-img/p1_chem_s2_q06_1.png}{51.3}{0}
\qfigdef{p1_chem_s3_q08_2}{Score Advanced - Test 38 (13 Apr 2025)-img/p1_chem_s3_q08_2.png}{69.4}{0}
\qfigdef{p1_chem_s3_q08_3}{Score Advanced - Test 38 (13 Apr 2025)-img/p1_chem_s3_q08_3.png}{69.9}{0}
\qfigdef{p1_chem_s3_q08_4}{Score Advanced - Test 38 (13 Apr 2025)-img/p1_chem_s3_q08_4.png}{68.5}{0}
\qfigdef{p1_chem_s3_q08_5}{Score Advanced - Test 38 (13 Apr 2025)-img/p1_chem_s3_q08_5.png}{70.7}{0}
\qfigdef{p1_chem_s3_q08_6}{Score Advanced - Test 38 (13 Apr 2025)-img/p1_chem_s3_q08_6.png}{70.8}{0}
\qfigdef{p1_chem_s3_q10_1}{Score Advanced - Test 38 (13 Apr 2025)-img/p1_chem_s3_q10_1.png}{124.0}{0}
\qfigdef{p1_chem_s3_q10_5}{Score Advanced - Test 38 (13 Apr 2025)-img/p1_chem_s3_q10_5.png}{47.8}{0}
\qfigdef{p1_chem_s3_q10_2}{Score Advanced - Test 38 (13 Apr 2025)-img/p1_chem_s3_q10_2.png}{146.8}{0}
\qfigdef{p1_chem_s3_q10_6}{Score Advanced - Test 38 (13 Apr 2025)-img/p1_chem_s3_q10_6.png}{49.6}{0}
\qfigdef{p1_chem_s3_q10_3}{Score Advanced - Test 38 (13 Apr 2025)-img/p1_chem_s3_q10_3.png}{152.3}{0}
\qfigdef{p1_chem_s3_q10_7}{Score Advanced - Test 38 (13 Apr 2025)-img/p1_chem_s3_q10_7.png}{49.6}{0}
\qfigdef{p1_chem_s3_q10_4}{Score Advanced - Test 38 (13 Apr 2025)-img/p1_chem_s3_q10_4.png}{151.8}{0}
\qfigdef{p1_chem_s3_q10_8}{Score Advanced - Test 38 (13 Apr 2025)-img/p1_chem_s3_q10_8.png}{48.0}{0}
\qfigdef{p1_chem_s3_q10_9}{Score Advanced - Test 38 (13 Apr 2025)-img/p1_chem_s3_q10_9.png}{49.4}{0}
\qfigdef{p1_chem_s3_q11_1}{Score Advanced - Test 38 (13 Apr 2025)-img/p1_chem_s3_q11_1.png}{110.7}{0}
\qfigdef{p1_chem_s3_q11_5}{Score Advanced - Test 38 (13 Apr 2025)-img/p1_chem_s3_q11_5.png}{54.7}{0}
\qfigdef{p1_chem_s3_q11_2}{Score Advanced - Test 38 (13 Apr 2025)-img/p1_chem_s3_q11_2.png}{132.5}{0}
\qfigdef{p1_chem_s3_q11_6}{Score Advanced - Test 38 (13 Apr 2025)-img/p1_chem_s3_q11_6.png}{43.4}{0}
\qfigdef{p1_chem_s3_q11_3}{Score Advanced - Test 38 (13 Apr 2025)-img/p1_chem_s3_q11_3.png}{87.5}{0}
\qfigdef{p1_chem_s3_q11_7}{Score Advanced - Test 38 (13 Apr 2025)-img/p1_chem_s3_q11_7.png}{60.4}{0}
\qfigdef{p1_chem_s3_q11_4}{Score Advanced - Test 38 (13 Apr 2025)-img/p1_chem_s3_q11_4.png}{117.7}{0}
\qfigdef{p1_chem_s3_q11_8}{Score Advanced - Test 38 (13 Apr 2025)-img/p1_chem_s3_q11_8.png}{40.2}{0}
\qfigdef{p1_chem_s3_q11_9}{Score Advanced - Test 38 (13 Apr 2025)-img/p1_chem_s3_q11_9.png}{56.5}{0}
\qfigdef{p1_chem_s4_q01_1}{Score Advanced - Test 38 (13 Apr 2025)-img/p1_chem_s4_q01_1.png}{117.0}{0}
\qfigdef{p1_chem_s4_q02_1}{Score Advanced - Test 38 (13 Apr 2025)-img/p1_chem_s4_q02_1.png}{183.2}{0}
\qfigdef{p1_chem_s4_q05_1}{Score Advanced - Test 38 (13 Apr 2025)-img/p1_chem_s4_q05_1.png}{126.7}{0}
\qfigdef{p2_phy_s1_q01_1}{Score Advanced - Test 38 (13 Apr 2025)-img/p2_phy_s1_q01_1.png}{123.0}{0}
\qfigdef{p2_phy_s1_q01_A}{Score Advanced - Test 38 (13 Apr 2025)-img/p2_phy_s1_q01_A.png}{119.4}{0}
\qfigdef{p2_phy_s1_q01_B}{Score Advanced - Test 38 (13 Apr 2025)-img/p2_phy_s1_q01_B.png}{117.0}{0}
\qfigdef{p2_phy_s1_q01_C}{Score Advanced - Test 38 (13 Apr 2025)-img/p2_phy_s1_q01_C.png}{119.6}{0}
\qfigdef{p2_phy_s1_q01_D}{Score Advanced - Test 38 (13 Apr 2025)-img/p2_phy_s1_q01_D.png}{118.4}{0}
\qfigdef{p2_phy_s1_q04_1}{Score Advanced - Test 38 (13 Apr 2025)-img/p2_phy_s1_q04_1.png}{195.5}{0}
\qfigdef{p2_phy_s2_q05_1}{Score Advanced - Test 38 (13 Apr 2025)-img/p2_phy_s2_q05_1.png}{57.2}{0}
\qfigdef{p2_phy_s2_q07_1}{Score Advanced - Test 38 (13 Apr 2025)-img/p2_phy_s2_q07_1.png}{151.0}{0}
\qfigdef{p2_phy_s3_para1_2_1}{Score Advanced - Test 38 (13 Apr 2025)-img/p2_phy_s3_para1_2_1.png}{106.3}{0}
\qfigdef{p2_phy_s3_para3_4_1}{Score Advanced - Test 38 (13 Apr 2025)-img/p2_phy_s3_para3_4_1.png}{172.4}{0}
\qfigdef{p2_phy_s4_q07_1}{Score Advanced - Test 38 (13 Apr 2025)-img/p2_phy_s4_q07_1.png}{62.6}{0}
\qfigdef{p2_phy_s4_q08_1}{Score Advanced - Test 38 (13 Apr 2025)-img/p2_phy_s4_q08_1.png}{113.8}{0}
\qfigdef{p2_phy_s4_q09_1}{Score Advanced - Test 38 (13 Apr 2025)-img/p2_phy_s4_q09_1.png}{104.9}{0}
\qfigdef{p2_chem_s1_q03_1}{Score Advanced - Test 38 (13 Apr 2025)-img/p2_chem_s1_q03_1.png}{194.0}{0}
\qfigdef{p2_chem_s1_q03_A}{Score Advanced - Test 38 (13 Apr 2025)-img/p2_chem_s1_q03_A.png}{50.8}{0}
\qfigdef{p2_chem_s1_q03_B}{Score Advanced - Test 38 (13 Apr 2025)-img/p2_chem_s1_q03_B.png}{57.5}{0}
\qfigdef{p2_chem_s1_q03_C}{Score Advanced - Test 38 (13 Apr 2025)-img/p2_chem_s1_q03_C.png}{23.8}{0}
\qfigdef{p2_chem_s1_q03_D}{Score Advanced - Test 38 (13 Apr 2025)-img/p2_chem_s1_q03_D.png}{52.8}{0}
\qfigdef{p2_chem_s2_q06_1}{Score Advanced - Test 38 (13 Apr 2025)-img/p2_chem_s2_q06_1.png}{235.2}{0}
\qfigdef{p2_chem_s2_q07_1}{Score Advanced - Test 38 (13 Apr 2025)-img/p2_chem_s2_q07_1.png}{80.8}{0}
\qfigdef{p2_chem_s4_q07_1}{Score Advanced - Test 38 (13 Apr 2025)-img/p2_chem_s4_q07_1.png}{146.6}{0}
\raggedcolumns
\makeatletter\@beginparpenalty=10000\makeatother
\newcolumntype{L}[1]{>{\raggedright\arraybackslash}p{#1}}
\begin{document}
\begin{center}{\Large\bfseries Score Advanced --- Test 38 (13 Apr 2025)}\end{center}
\begin{multicols}{2}
\raggedcolumns
\papersec{1}
\partsec{1}{PHYSICS}
\sect{SECTION-I(i)}
\qn{1} A dimensionless quantity is constructed in terms of electronic charge e, permittivity of free space $\varepsilon_0$, Planck's constant h, and speed of light c. If the dimensionless quantity is written as e$^\alpha\varepsilon_0^\beta$h$^\gamma$c$^\delta$ and \textit{n} is a non-zero integer, then ($\alpha$, $\beta$, $\gamma$, $\delta$) is given by
\begin{qoptions}
\opt{A}{(\textit n, $-\textit n$, $-\textit n$, $-2\textit n$)}
\opt{B}{(2\textit n, $-\textit n$, $-2\textit n$, $-2\textit n$)}
\opt{C}{(2\textit n, $-\textit n$, $-\textit n$, $-\textit n$)}
\opt{D}{(\textit n, $-\textit n$, $-2\textit n$, $-\textit n$)}
\end{qoptions}
\qn{2} An infinitely long wire, located on the z-axis, carries a current \textit I along the +z-direction and produces the magnetic field $\vec{\mathrm B}$. The magnitude of the line integral $\displaystyle\int\vec{\mathrm B}\cdot\overrightarrow{\mathrm{dl}}$ along a straight line from the point ($-\sqrt3$\textit a, \textit a, 0) to ($\sqrt3$\textit a, \textit a, 0) is given by : [$\mu_0$ is the magnetic permeability of free space.]
\begin{qoptions}
\opt{A}{7$\mu_0$\textit I /24}
\opt{B}{7$\mu_0$\textit I /12}
\opt{C}{$\mu_0$\textit I /3}
\opt{D}{$\mu_0$\textit I /6}
\end{qoptions}
\qn{3} Two beads, each with charge \textit q and mass \textit m, are on a horizontal, frictionless, non-conducting, circular hoop of radius \textit R. One of the beads is glued to the hoop at some point, while the other one performs small oscillations about its equilibrium position along the hoop. The square of the angular frequency of the small oscillations is given by [$\varepsilon_0$ is the permittivity of free space]
\begin{qoptions}
\opt{A}{\textit q$^2$ / (8$\pi\varepsilon_0$\textit R$^3$\textit m)}
\opt{B}{\textit q$^2$ / (16$\pi\varepsilon_0$\textit R$^3$\textit m)}
\opt{C}{\textit q$^2$ / (32$\pi\varepsilon_0$\textit R$^3$\textit m)}
\opt{D}{\textit q$^2$ / (4$\pi\varepsilon_0$\textit R$^3$\textit m)}
\end{qoptions}
\qn{4} A block of mass 2.5 kg moves along the x-direction subject to the force \textit F = ($-$10\textit x + 5) N, with the value of \textit x in metre. At time \textit t = 0 s, it is at rest at position \textit x = 1 m. The position and momentum of the block at \textit t = ($\pi$/4) s are
\begin{qoptions}
\opt{A}{$-$0.5 m, 2.5 kg m/s}
\opt{B}{0.5 m, 0 kg m/s}
\opt{C}{0.5 m, $-$2.5 kg m/s}
\opt{D}{$-$1 m, 5 kg m/s}
\end{qoptions}
\sect{SECTION-I(ii)}
\qn{5} A particle of mass \textit m is moving in a circular orbit under the influence of the central force F(\textit r) = $-$\textit{kr}, corresponding to the potential energy V(\textit r) = \textit{kr}$^2$ / 2, where \textit k is a positive force constant and \textit r is the radial distance from the origin. According to the Bohr's quantization rule, the angular momentum of the particle is given by L = \textit{n}$\hbar$, where $\hbar$ = \textit h/(2$\pi$), \textit h is the Planck's constant, and \textit n a positive integer. If \textit v and \textit E are the speed and total energy of the particle, respectively, then which of the following expression(s) is(are) correct ?
\begin{qoptions}
\opt{A}{$\dfrac{\mathrm L}{\mathrm{mr}^2}=\sqrt{\dfrac{\mathrm k}{\mathrm m}}$}
\opt{B}{E = $\dfrac{\mathrm n\hbar}2\sqrt{\dfrac{\mathrm k}{\mathrm m}}$}
\opt{C}{v$^2$ = n$\hbar\sqrt{\dfrac{\mathrm k}{\mathrm m^3}}$}
\opt{D}{r$^2$ = n$\hbar\sqrt{\dfrac1{\mathrm{mk}}}$}
\end{qoptions}
\qn{6} Two uniform string of mass per unit length $\mu$ and 4$\mu$, and length \textit L and 2\textit L, respectively, are joined at point O, and tied at two fixed ends P and Q, as shown in the figure. The strings are under a uniform tension T. If we define the frequency v$_0$ = $\dfrac1{2L}\sqrt{\dfrac{\mathrm T}{\mu}}$, which of the following statement(s) is (are) correct ?\par
\qfig{p1_phy_s2_q06_1}
\begin{qoptions}
\opt{A}{No vibrational mode with an antinode at O is possible for the composite string.}
\opt{B}{When the composite string vibrates at the minimum frequency with a node at O, it has 6 nodes, including the end nodes.}
\opt{C}{With a node at O, the minimum frequency of vibration of the composite string is \textit v$_0$.}
\opt{D}{With an antinode at O, the minimum frequency of vibration of the composite string is 2\textit v$_0$.}
\end{qoptions}
\qn{7} A glass beaker has a solid, plano-convex base of refractive index 1.60, as shown in the figure. The radius of curvature of the convex surface (SPU) is 9 cm, while the planar surface (STU) acts as a mirror. This beaker is filled with a liquid of refractive index \textit n up to the level QPR. If the image of a point object O at a height of \textit h (OT in the figure) is formed onto itself, then, which of the following option(s) is(are) correct ?\par
\qfig{p1_phy_s2_q07_1}
\begin{qoptions}
\opt{A}{For \textit n = 1.42, \textit h = 40 cm.}
\opt{B}{For \textit n = 1.35, \textit h = 52 cm.}
\opt{C}{For \textit n = 1.45, \textit h = 60 cm.}
\opt{D}{For \textit n = 1.48, \textit h = 75 cm.}
\end{qoptions}
\sect{SECTION-I(iii)}
\qn{8} One mole of a monoatomic ideal gas undergoes the cyclic process J $\to$ K $\to$ L $\to$ M $\to$ J, as shown in the P-T diagram.\par
\qfig{p1_phy_s3_q08_1}
Match the quantities mentioned in List-I with their values in List-II and choose the correct option. [\textit R is the gas constant.]\par
{\footnotesize\setlength{\tabcolsep}{1.5pt}\noindent\begin{tabular}{|c|L{4.2cm}|c|L{2.4cm}|}
\hline
 & \multicolumn{1}{>{\centering\arraybackslash}p{4.2cm}|}{List-I} & & \multicolumn{1}{>{\centering\arraybackslash}p{2.4cm}|}{List-II} \\ \hline
(P) & Work done in the complete cyclic process & (1) & 0 \\ \hline
(Q) & Change in the internal energy of the gas in the process JK & (2) & $-3RT_0\ln2$ \\ \hline
(R) & Heat given to the gas in the process KL & (3) & 3\textit{RT}$_0$ \\ \hline
(S) & Change in the internal energy of the gas in the process MJ & (4) & \textit{RT}$_0$ $-$ 4RT$_0$ ln 2 \\ \hline
 & & (5) & $-2RT_0\ln2$ \\ \hline
\end{tabular}}\par
\begin{qoptions}
\opt{A}{P $\to$ 1;Q $\to$ 3;R $\to$ 5;S $\to$ 4}
\opt{B}{P $\to$ 5;Q $\to$ 3;R $\to$ 2;S $\to$ 1}
\opt{C}{P $\to$ 5;Q $\to$ 4;R $\to$ 1;S $\to$ 1}
\opt{D}{P $\to$ 2;Q $\to$ 5;R $\to$ 3;S $\to$ 4}
\end{qoptions}
\qn{9} Four identical thin, square metal sheets, \textit S$_1$, \textit S$_2$, \textit S$_3$ and \textit S$_4$, each of side \textit a are kept parallel to each other with equal distance \textit d ($\ll$ \textit a) between them, as shown in the figure. Let \textit C$_0$ = $\varepsilon_0$\textit a$^2$/d, where $\varepsilon_0$ is the permittivity of free space.\par
\qfig{p1_phy_s3_q09_1}
Match the quantities mentioned in List-I with their values in List-II and choose the correct option.\par
{\footnotesize\setlength{\tabcolsep}{1.5pt}\noindent\begin{tabular}{|c|L{4.6cm}|c|L{1.4cm}|}
\hline
 & \multicolumn{1}{>{\centering\arraybackslash}p{4.6cm}|}{List-I} & & \multicolumn{1}{>{\centering\arraybackslash}p{1.4cm}|}{List-II} \\ \hline
(P) & The capacitance between \textit S$_1$ and \textit S$_4$, with \textit S$_2$ shorted to \textit S$_3$, is & (1) & 3\textit C$_0$ \\ \hline
(Q) & The capacitance between \textit S$_1$ and \textit S$_2$, with \textit S$_3$ shorted to \textit S$_1$, and \textit S$_2$ shorted to \textit S$_4$, is & (2) & \textit C$_0$/2 \\ \hline
(R) & The capacitance between \textit S$_1$ and \textit S$_4$, with \textit S$_2$ and \textit S$_3$ not connected, is & (3) & \textit C$_0$/3 \\ \hline
(S) & The capacitance between \textit S$_1$ and \textit S$_3$, with \textit S$_2$ shorted to \textit S$_4$, is & (4) & 2\textit C$_0$/3 \\ \hline
 & & (5) & 2\textit C$_0$ \\ \hline
\end{tabular}}\par
\begin{qoptions}
\opt{A}{P $\to$ 2;Q $\to$ 1;R $\to$ 3;S $\to$ 4}
\opt{B}{P $\to$ 4;Q $\to$ 2;R $\to$ 3;S $\to$ 1}
\opt{C}{P $\to$ 2;Q $\to$ 4;R $\to$ 1;S $\to$ 4}
\opt{D}{P $\to$ 4;Q $\to$ 3;R $\to$ 2;S $\to$ 1}
\end{qoptions}
\qn{10} A light ray is incident on the surface of a sphere of refractive index n at an angle of incidence $\theta_0$. The ray partially refracts into the sphere with angle of refraction $\phi_0$ and then partly reflects from the back surface. The reflected ray then emerges out of the sphere after a partial refraction. The total angle of deviation of the emergent ray with respect to the incident ray is $\alpha$. Match the quantities mentioned in List-I with their values in List-II and choose the correct option.\par
{\footnotesize\setlength{\tabcolsep}{1.5pt}\noindent\begin{tabular}{|c|L{5.0cm}|c|L{2.0cm}|}
\hline
 & \multicolumn{1}{>{\centering\arraybackslash}p{5.0cm}|}{List-I} & & \multicolumn{1}{>{\centering\arraybackslash}p{2.0cm}|}{List-II} \\ \hline
(P) & If \textit n = 2 and $\alpha$ = 180$^\circ$, then all the possible values of $\theta_0$ will be & (1) & 60$^\circ$ and 0$^\circ$ \\ \hline
(Q) & If n = $\sqrt3$ and $\alpha$ = 180$^\circ$, then all the possible values of $\theta_0$ will be & (2) & 150$^\circ$ \\ \hline
(R) & If n = $\sqrt3$ and $\alpha$ = 180$^\circ$, then all the possible values of $\phi_0$ will be & (3) & 0$^\circ$ \\ \hline
(S) & If n = $\sqrt2$ and $\theta_0$ = 45$^\circ$, then all the possible values of $\alpha$ will be & (4) & 30$^\circ$ and 0$^\circ$ \\ \hline
 & & (5) & 45$^\circ$ and 0$^\circ$ \\ \hline
\end{tabular}}\par
\begin{qoptions}
\opt{A}{P $\to$ 3;Q $\to$ 2;R $\to$ 4;S $\to$ 3}
\opt{B}{P $\to$ 5;Q $\to$ 1;R $\to$ 2;S $\to$ 4}
\opt{C}{P $\to$ 3;Q $\to$ 1;R $\to$ 4;S $\to$ 2}
\opt{D}{P $\to$ 3;Q $\to$ 1;R $\to$ 2;S $\to$ 5}
\end{qoptions}
\qn{11} The circuit shown in the figure contains an inductor \textit L, a capacitor \textit C$_0$, a resistor \textit R$_0$ and an ideal battery. The circuit also contains two keys K$_1$ and K$_2$. Initially, both the keys are open and there is no charge on the capacitor. At an instant, key K$_1$ is closed and immediately after this the current in \textit R$_0$ is found to be \textit I$_1$. After a long time, the current attains a steady state value \textit I$_2$. Thereafter, K$_2$ is closed and simultaneously K$_1$ is opened and the voltage across \textit C$_0$ oscillates with amplitude \textit V$_0$ and angular frequency $\omega_0$.\par
\qfig{p1_phy_s3_q11_1}
Match the quantities mentioned in List-I with their values in List-II and choose the correct option.\par
{\footnotesize\setlength{\tabcolsep}{1.5pt}\noindent\begin{tabular}{|c|L{4.0cm}|c|L{1.2cm}|}
\hline
 & \multicolumn{1}{>{\centering\arraybackslash}p{4.0cm}|}{List-I} & & \multicolumn{1}{>{\centering\arraybackslash}p{1.2cm}|}{List-II} \\ \hline
(P) & The value of $\omega_0$ in kilo-radians/s is & (1) & 0 \\ \hline
(Q) & The value of V$_0$ in Volt is & (2) & 2 \\ \hline
(R) & The value of \textit I$_1$ in Ampere is & (3) & 4 \\ \hline
(S) & The value of \textit I$_2$ in Ampere is & (4) & 20 \\ \hline
 & & (5) & 200 \\ \hline
\end{tabular}}\par
\begin{qoptions}
\opt{A}{P $\to$ 2;Q $\to$ 5;R $\to$ 1;S $\to$ 3}
\opt{B}{P $\to$ 2;Q $\to$ 1;R $\to$ 4;S $\to$ 3}
\opt{C}{P $\to$ 1;Q $\to$ 3;R $\to$ 2;S $\to$ 5}
\opt{D}{P $\to$ 4;Q $\to$ 3;R $\to$ 2;S $\to$ 1}
\end{qoptions}
\sect{SECTION-II}
\qn{1} The specific heat capacity of a substance is temperature dependent and is given by the formula \textit C = \textit{kT}, where \textit k is a constant of suitable dimensions in SI units, and \textit T is the absolute temperature. If the heat required to raise the temperature of 1 kg of the substance from $-$23$^\circ$C to 27$^\circ$C is \textit{nk}, the value of $\dfrac{\mathrm n}{50}$ is \underline{\hspace{1.5cm}}.\par
[Given : 0 K = $-$273 $^\circ$C.]
\qn{2} A disc of mass \textit M and radius \textit R is free to rotate about its vertical axis as shown in the figure. A battery operated motor of negligible mass is fixed to this disc at a point on its circumference. Another disc of the mass $\dfrac{\mathrm M}2$ and radius \textit R / 2 is fixed to the motor's thin shaft. Initially, both the discs are at rest. The motor is switched on so that the smaller disc rotates at a uniform angular speed $\omega$ about its axis as seen from ground. If the angular speed at which the large disc rotates is $\omega$ / \textit n, then the value of \textit n is \underline{\hspace{1.5cm}}.\par
\qfig{p1_phy_s4_q02_1}
\qn{3} A point source S emits unpolarized light uniformly in all directions. At two points A and B, on the same line as S the ratio r = I$_\mathrm A$/ I$_\mathrm B$ of the intensities of light is 2. If a polaroids having 45$^\circ$ angle between their pass-axes is placed just before point A and just before point B respectively, then the new value of r will be \underline{\hspace{1.5cm}}.\par
\qfig{p1_phy_s4_q03_1}
\qn{4} A source (S) of sound has frequency 240 Hz. When the observer (O) and the source move away from each other at a speed \textit v with respect to the ground (as shown in Case 1 in the figure), the observer measures the frequency of the sound to be 200 Hz. However, when the observer and the source move towards each other at the same speed \textit v with respect to the ground (as shown in Case 2 in the figure), the observer measures the frequency of sound to be \textit n Hz. The value of \textit n is \underline{\hspace{1.5cm}}.\par
\qfig{p1_phy_s4_q04_1}
\qn{5} Two large, identical water tanks, 1 and 2, kept on the top of a building of height \textit H, are filled with water up to height \textit h in each tank. Both the tanks contain an identical hole of small radius on their sides, close to their bottom. A pipe of the same internal radius as that of the hole is connected to tank 2, and the pipe ends at the ground level. When the water flows from the tanks 1 and 2 through the holes, the times taken to empty the tanks are \textit t$_1$ and \textit t$_2$, respectively. If H = $\left(\dfrac{16}9\right)$\textit h, then t$_1$/t$_2$ is \underline{\hspace{1.5cm}}.
\qn{6} A thin uniform rod of length \textit L and certain mass is kept on a frictionless horizontal table with a massless string of length \textit L fixed to one end (top view is shown in the figure). The other end of the string is pivoted to a point O. If a horizontal impulse \textit P is imparted to the rod at a distance \textit x = \textit L/n from the mid-point of the rod (see figure), then the rod and string revolve together around the point O, with the rod remaining aligned with the string. In such a case, the value of \textit n is \underline{\hspace{1.5cm}}.\par
\qfig{p1_phy_s4_q06_1}
\raggedcolumns
\papersec{1}
\partsec{2}{CHEMISTRY}
\sect{SECTION-I(i)}
\qn{1} A closed vessel contains 10 g of an ideal gas \textbf X at 300 K, which exerts 2 atm pressure. At the same temperature, 80 g of another ideal gas \textbf Y is added to it and the pressure becomes 6 atm. The ratio of root mean square velocities of \textbf Y and \textbf X at 300 K is
\begin{qoptions}
\opt{A}{2$\sqrt2$ : $\sqrt3$}
\opt{B}{2$\sqrt2$ : 1}
\opt{C}{1 : 2}
\opt{D}{2 : 1}
\end{qoptions}
\qn{2} Aspartame, an artificial sweetener, is a dipeptide aspartyl phenylalanine methyl ester. The structure of aspartame is\par
\qfig{p1_chem_s1_q02_1}
\begin{qoptions}
\opt{A}{\qfigi{p1_chem_s1_q02_A}}
\opt{B}{\qfigi{p1_chem_s1_q02_B}}
\opt{C}{\qfigi{p1_chem_s1_q02_C}}
\opt{D}{\qfigi{p1_chem_s1_q02_D}}
\end{qoptions}
\qn{3} The species formed on fluorination of phosphorus pentachloride in a polar organic solvent are :
\begin{qoptions}
\opt{A}{[PF$_4$]$^+$ [PF$_6$]$^-$ and [PCl$_4$]$^+$ [PF$_6$]$^-$}
\opt{B}{PF$_3$ and PCl$_3$}
\opt{C}{[PCl$_4$]$^+$ [PCl$_4$F$_2$]$^-$ and [PCl$_4$]$^+$ [PF$_6$]$^-$}
\opt{D}{PF$_5$ and PCl$_3$}
\end{qoptions}
\qn{4} The option in which at least three molecules follow Octet Rule is
\begin{qoptions}
\opt{A}{HCN, C$_2$H$_2$ , ClO$_2$ , H$_3$PO$_4$}
\opt{B}{NO$_2$ , O$_3$ , NH$_3$ , NO}
\opt{C}{BF$_3$ , ClO$_3$ , CH$_4$ , HClO$_4$}
\opt{D}{XeF$_2$ , C$_3$H$_4$ , O$_3$ , C$_6$H$_6$}
\end{qoptions}
\sect{SECTION-I(ii)}
\qn{5} Among the following, the correct statement(s) for electrons in an atom is(are)
\begin{qoptions}
\opt{A}{According to Bohr's model, the magnitude of velocity of electrons increases with decrease in values of n.}
\opt{B}{According to Bohr's model, the most negative energy value for an electron is given by n = 1, which corresponds to the most stable orbit.}
\opt{C}{Uncertainty principle rules out the existence of definite paths for electrons.}
\opt{D}{The energy of an electron in 2\textit s orbital of an atom is higher than the energy of an electron that is infinitely far away from the nucleus.}
\end{qoptions}
\qn{6} Reaction of \textit{iso}-propylbenzene with O$_2$ followed by the treatment with H$_3$O$^+$ forms phenol and a by-product \textbf P. Reaction of \textbf P with 3 equivalents of Cl$_2$ gives compound \textbf Q. Treatment of \textbf Q with Ca(OH)$_2$ produces compound \textbf R and calcium salt \textbf S.\par
The correct statement(s) regarding \textbf P, \textbf Q, \textbf R and \textbf S is(are)
\begin{qoptions}
\opt{A}{Reaction of \textbf P with \textbf R in the presence of KOH followed by acidification gives \qfigi{p1_chem_s2_q06_1}}
\opt{B}{Reaction of \textbf R with O$_2$ in the presence of light gives phosgene gas}
\opt{C}{\textbf Q react with CO$_2$ in presence of NaOH followed by H$_3$O$^\oplus$ gives salicylic acid.}
\opt{D}{\textbf S on heating gives \textbf P}
\end{qoptions}
\qn{7} The compound (s) having peroxide linkage is(are) :
\begin{qoptions}
\opt{A}{H$_3$PO$_5$}
\opt{B}{H$_2$S$_2$O$_8$}
\opt{C}{H$_4$P$_2$O$_8$}
\opt{D}{H$_2$S$_2$O$_5$}
\end{qoptions}
\sect{SECTION-I(iii)}
\qn{8} In a conductometric titration, small volume of titrant of higher concentration is added stepwise to a larger volume of titrate of much lower concentration, and the conductance is measured after each addition.\par
The limiting ionic conductivity ($\Lambda_0$) values (in mS m$^2$ mol$^{-1}$) for different ions in aqueous solutions are given below :\par
{\footnotesize\setlength{\tabcolsep}{2pt}\noindent\begin{tabular}{|c|c|c|c|c|c|c|c|c|c|}
\hline
Ions & Ag$^+$ & K$^+$ & Na$^+$ & H$^+$ & NO$_3^-$ & Cl$^-$ & SO$_4^{2-}$ & OH$^-$ & CH$_3$COO$^-$ \\ \hline
$\Lambda_0$ & 6.2 & 7.4 & 5.0 & 35.0 & 7.2 & 7.6 & 16.0 & 19.9 & 4.1 \\ \hline
\end{tabular}}\par
For different combinations of titrates and titrants given in \textbf{List-I}, the graphs of `conductance' versus `volume of titrant' are given in \textbf{List-II}.\par
Match each entry in \textbf{List-I} with the appropriate entry in \textbf{List-II} and choose the correct option.\par
{\footnotesize\setlength{\tabcolsep}{1.5pt}\noindent\begin{tabular}{|c|L{3.2cm}|c|c|}
\hline
 & \multicolumn{1}{>{\centering\arraybackslash}p{3.2cm}|}{List-I} & & \multicolumn{1}{c|}{List-II} \\ \hline
(P) & Titrate : AgNO$_3$ \newline Titrant : KCl & (1) & \qfigi{p1_chem_s3_q08_2} \\ \hline
(Q) & Titrate : KCl \newline Titrant : AgNO$_3$ & (2) & \qfigi{p1_chem_s3_q08_3} \\ \hline
(R) & Titrate : NaOH \newline Titrant : HCl & (3) & \qfigi{p1_chem_s3_q08_4} \\ \hline
(S) & Titrate : CH$_3$COOH \newline Titrant : NaOH & (4) & \qfigi{p1_chem_s3_q08_5} \\ \hline
 & & (5) & \qfigi{p1_chem_s3_q08_6} \\ \hline
\end{tabular}}\par
\begin{qoptions}
\opt{A}{P $\to$ 4;Q $\to$ 3;R $\to$ 2;S $\to$ 5}
\opt{B}{P $\to$ 2;Q $\to$ 4;R $\to$ 3;S $\to$ 1}
\opt{C}{P $\to$ 3;Q $\to$ 4;R $\to$ 2;S $\to$ 5}
\opt{D}{P $\to$ 4;Q $\to$ 3;R $\to$ 2;S $\to$ 1}
\end{qoptions}
\qn{9} Based on \textbf{VSEPR} model, match the xenon compounds given in \textbf{List-I} with the corresponding geometries and the number of lone pairs on xenon given in \textbf{List-II} and choose the correct option.\par
{\footnotesize\setlength{\tabcolsep}{1.5pt}\noindent\begin{tabular}{|c|L{1.5cm}|c|L{4.6cm}|}
\hline
 & \multicolumn{1}{>{\centering\arraybackslash}p{1.5cm}|}{List-I} & & \multicolumn{1}{>{\centering\arraybackslash}p{4.6cm}|}{List-II} \\ \hline
(P) & XeO$_3$ & (1) & Trigonal bipyramidal and two lone pair of electrons \\ \hline
(Q) & XeF$_4$ & (2) & Tetrahedral and one lone pair of electrons \\ \hline
(R) & XeO$_2$F$_2$ & (3) & Octahedral and two lone pair of electrons \\ \hline
(S) & XeOF$_4$ & (4) & Trigonal bipyramidal and one lone pair of electrons \\ \hline
 & & (5) & Octahedral and one lone pair of electron \\ \hline
\end{tabular}}\par
\begin{qoptions}
\opt{A}{P $\to$ 5;Q $\to$ 2;R $\to$ 3;S $\to$ 1}
\opt{B}{P $\to$ 2;Q $\to$ 3;R $\to$ 4;S $\to$ 1}
\opt{C}{P $\to$ 4;Q $\to$ 3;R $\to$ 2;S $\to$ 1}
\opt{D}{P $\to$ 2;Q $\to$ 3;R $\to$ 4;S $\to$ 5}
\end{qoptions}
\qn{10} \textbf{List-I} contains various reaction sequences and \textbf{List-II} contains the possible products. Match each entry in \textbf{List-I} with the appropriate entry in \textbf{List-II} and choose the correct option.\par
{\footnotesize\setlength{\tabcolsep}{1.5pt}\noindent\begin{tabular}{|c|c|c|c|}
\hline
 & \multicolumn{1}{c|}{List-I} & & \multicolumn{1}{c|}{List-II} \\ \hline
(P) & \qfigi{p1_chem_s3_q10_1} & (1) & \qfigi{p1_chem_s3_q10_5} \\ \hline
(Q) & \qfigi{p1_chem_s3_q10_2} & (2) & \qfigi{p1_chem_s3_q10_6} \\ \hline
(R) & \qfigi{p1_chem_s3_q10_3} & (3) & \qfigi{p1_chem_s3_q10_7} \\ \hline
(S) & \qfigi{p1_chem_s3_q10_4} & (4) & \qfigi{p1_chem_s3_q10_8} \\ \hline
 & & (5) & \qfigi{p1_chem_s3_q10_9} \\ \hline
\end{tabular}}\par
\begin{qoptions}
\opt{A}{P $\to$ 3;Q $\to$ 2;R $\to$ 1;S $\to$ 4}
\opt{B}{P $\to$ 3;Q $\to$ 5;R $\to$ 4;S $\to$ 1}
\opt{C}{P $\to$ 3;Q $\to$ 2;R $\to$ 4;S $\to$ 1}
\opt{D}{P $\to$ 2;Q $\to$ 5;R $\to$ 4;S $\to$ 1}
\end{qoptions}
\qn{11} \textbf{List-I} contains various reaction sequences and \textbf{List-II} contains different phenolic compounds. Match each entry in \textbf{List-I} with the appropriate entry in \textbf{List-II} and choose the correct option.\par
{\footnotesize\setlength{\tabcolsep}{1.5pt}\noindent\begin{tabular}{|c|c|c|c|}
\hline
 & \multicolumn{1}{c|}{List-I} & & \multicolumn{1}{c|}{List-II} \\ \hline
(P) & \qfigi{p1_chem_s3_q11_1} & (1) & \qfigi{p1_chem_s3_q11_5} \\ \hline
(Q) & \qfigi{p1_chem_s3_q11_2} & (2) & \qfigi{p1_chem_s3_q11_6} \\ \hline
(R) & \qfigi{p1_chem_s3_q11_3} & (3) & \qfigi{p1_chem_s3_q11_7} \\ \hline
(S) & \qfigi{p1_chem_s3_q11_4} & (4) & \qfigi{p1_chem_s3_q11_8} \\ \hline
 & & (5) & \qfigi{p1_chem_s3_q11_9} \\ \hline
\end{tabular}}\par
\begin{qoptions}
\opt{A}{P $\to$ 2;Q $\to$ 5;R $\to$ 4;S $\to$ 3}
\opt{B}{P $\to$ 5;Q $\to$ 3;R $\to$ 4;S $\to$ 1}
\opt{C}{P $\to$ 2;Q $\to$ 5;R $\to$ 3;S $\to$ 1}
\opt{D}{P $\to$ 5;Q $\to$ 2;R $\to$ 3;S $\to$ 4}
\end{qoptions}
\sect{SECTION-II}
\qn{1} Consider the following volume-temperature (V $-$ T) diagram for the expansion of 5 moles of an ideal monoatomic gas.\par
\qfig{p1_chem_s4_q01_1}
Considering only P-V work is involved, the total change in internal energy (in Joule) for the transformation of state in the sequence \textbf X $\to$ \textbf Y $\to$ \textbf Z is \underline{\hspace{1.5cm}}.\par
[Use the given data: Molar heat capacity of the gas for the given temperature range, C$_{\mathrm{V,\,m}}$ = 12 J K$^{-1}$ mol$^{-1}$ and gas constant, R = 8.3 J K$^{-1}$ mol$^{-1}$]
\qn{2} Consider the following reaction,\par
2H$_2$(g) + 2NO(g) $\longrightarrow$ N$_2$(g) + 2H$_2$O(g)\par
which follows the mechanism given below:\par
\qfig{p1_chem_s4_q02_1}
N$_2$O$_2$(g) + H$_2$(g) $\xrightarrow{k_2}$ N$_2$O(g) + H$_2$O(g) \quad (slow reaction)\par
N$_2$O(g) + H$_2$(g) $\xrightarrow{k_3}$ N$_2$(g) + H$_2$O(g) \quad (fast reaction)\par
The order with respect to H$_2$ is \underline{\hspace{1.5cm}}.
\qn{3} Complete reaction of acetaldehyde with excess formaldehyde, upon heating with conc. NaOH solution, gives \textbf P and \textbf Q. Compound \textbf P does not give Tollens' test, whereas \textbf Q on acidification gives positive Tollens' test. Treatment of \textbf P with excess cyclohexanone in the presence of catalytic amount of \textit p-toluenesulfonic acid (PTSA) gives product \textbf R.\par
Total number of methylene groups (-CH$_2$-) in \textbf R is \underline{\hspace{1.5cm}}.
\qn{4} Among [V(CO)$_6$]$^-$, Cr(CO)$_5$, Cu(CO)$_3$, [Mn(CO)$_6$]$^+$, Fe(CO)$_5$, [Co(CO)$_3$]$^{2-}$, [Cr(CO)$_4$]$^{4-}$, [Ti(CO)$_6$]$^{-2}$, [Co(CO)$_4$]$^{-1}$ and Ir(CO)$_3$, the total number of species isoelectronic with Ni(CO)$_4$ is \underline{\hspace{1.5cm}}.\par
[Given, atomic number : V = 23, Cr = 24, Mn = 25, Fe = 26, Co = 27, Ni = 28, Cu = 29, Ir = 77]
\qn{5} In the following reaction sequence, the major product \textbf P is formed.\par
\qfig{p1_chem_s4_q05_1}
1 mole of \textbf P react with $\dfrac13$ mole of glycerol in presence of an acid as catalyst gives product \textbf Q. Reaction of \textbf Q with excess NaOH followed by the treatment with CaCl$_2$ yields \textbf R, quantitatively.\par
Starting with $\dfrac13$ mole of \textbf Q, the amount of \textbf R produced in gm is \underline{\hspace{1.5cm}}\par
[Given, atomic weight: H = 1, C = 12, N = 14, O = 16, Na = 23, Cl = 35, Ca = 40]
\qn{6} Among the following complexes, the total number of diamagnetic species is \underline{\hspace{1.5cm}}.\par
[Mn(NH$_3$)$_6$]$^{3+}$, [MnCl$_6$]$^{3-}$, [Ni(CO)$_4$], [CoF$_6$]$^{3-}$, [Fe(NH$_3$)$_6$]$^{3+}$, O$_3$ , BaO$_2$ , [Pt(NH$_3$)$_2$Br$_2$], [Co(en)$_3$]Cl$_2$ and [Fe(CN)$_6$]$^{4-}$\par
[Given, atomic number : Mn = 25, Fe = 26, Co = 27;\par
en = H$_2$NCH$_2$CH$_2$NH$_2$]
\raggedcolumns
\papersec{1}
\partsec{3}{MATHEMATICS}
\sect{SECTION-I(i)}
\qn{1} Let f(x) is a continuous and differentiable on (0, $\infty$) such that f(1) = 2 and $\displaystyle\lim_{t\to x}\dfrac{t^3f(x)-x^3f(t)}{t^2-x^2}=2$ for each x $>$ 0. Then for all x $>$ 0, f(x) is equal to
\begin{qoptions}
\opt{A}{f (x) = x$^3$ + $\dfrac1x$}
\opt{B}{f (x) = x$^3$ + $\dfrac1{x^2}$}
\opt{C}{f (x) = 3x$^3$ $-$ $\dfrac1x$}
\opt{D}{f (x) = 3x$^3$ $-$ $\dfrac1{x^2}$}
\end{qoptions}
\qn{2} Let $\dfrac\pi2<$ x $<\pi$ and tanx = $\dfrac{-24}7$, then value of $\sin\dfrac{13x}2(\sin7x-\cos7x)+\cos\dfrac{13x}2(\sin7x+\cos7x)$ is equal to
\begin{qoptions}
\opt{A}{$-\dfrac75$}
\opt{B}{$\dfrac75$}
\opt{C}{$\dfrac43$}
\opt{D}{$-\dfrac43$}
\end{qoptions}
\qn{3} A student appears for a quiz consisting of only, objective type questions, each question having 4 options of which exactly one option is correct and he answers all the questions. The student knows the answers of some questions and guesses the answers for the remaining questions. Whenever the student knows the answer of a question, he gives the correct answer. Assume that the probability of the student giving the correct answer for a question, given that he has guessed it, is $\dfrac14$. Also assume that the probability of the answer for a question being guessed, given that the student's answer is correct, is $\dfrac16$. Then the probability that the student knows the answer of a randomly chosen question is :
\begin{qoptions}
\opt{A}{$\dfrac49$}
\opt{B}{$\dfrac5{12}$}
\opt{C}{$\dfrac57$}
\opt{D}{$\dfrac59$}
\end{qoptions}
\qn{4} Consider the ellipse $\dfrac{x^2}9+\dfrac{y^2}4=1$. Let S(p, q) be a point in the first quadrant such that $\dfrac{p^2}9+\dfrac{q^2}4>1$. Two tangents are drawn from S to the ellipse, of which one meets the ellipse at one end point of the minor axis and the other meets the ellipse at a point T in the fourth quadrant. Let R be the vertex of the ellipse with positive x-coordinate and O be the center of the ellipse. If the area of the triangle $\Delta$ORT is $\dfrac95$ then which of the following options is correct ?
\begin{qoptions}
\opt{A}{p = 3, q = 2}
\opt{B}{p = 6, q = 2}
\opt{C}{p = 3, q = 1}
\opt{D}{p = 6, q = 1}
\end{qoptions}
\sect{SECTION-I(ii)}
\qn{5} Let S = $\left\{\mathrm a+\mathrm b\sqrt3:\mathrm a,\mathrm b\in\mathrm Z\right\}$, T$_1$ = $\left\{\left(2-\sqrt3\right)^{\mathrm n}:\mathrm n\in\mathrm N\right\}$ and T$_2$ = $\left\{\left(2+\sqrt3\right)^{\mathrm n}:\mathrm n\in\mathrm N\right\}$, then which of the following is (are) true ?
\begin{qoptions}
\opt{A}{Z $\cup$ T$_1$ $\cup$ T$_2$ $\subset$ S}
\opt{B}{T$_2$ $\cap\left(0,\dfrac1{2025}\right)$ = $\phi$\par (where $\phi$ denotes null set)}
\opt{C}{T$_1$ $\cap\left(0,\dfrac1{100}\right)\neq\phi$}
\opt{D}{For any a, b $\in$ Z, $\cos\left(\pi\left(\mathrm a+\mathrm b\sqrt3\right)\right)+\mathrm i\sin\left(\pi\left(\mathrm a+\mathrm b\sqrt3\right)\right)\in$ Z if and only if b = 0, where i = $\sqrt{-1}$}
\end{qoptions}
\qn{6} Let R$^2$ denote R $\times$ R. Let S = \{(a, b, c) : a, b, c $\in$ R and ax$^2$ + 2bxy + cy$^2$ $<$ 0, $\forall$ (x, y) $\in$ R$^2$ $-$ \{(0, 0)\}\} then which of the following statements is(are) True ?
\begin{qoptions}
\opt{A}{($-$2, 3, $-$4) $\in$ S}
\opt{B}{If ($-$1, 2, c) $\in$ S, then c $<$ $-$4}
\opt{C}{For some given (a, b, c) $\in$ S, the system of equation $\begin{array}{l}\mathrm{ax+by}=1\\\mathrm{bx+2cy}=-1\end{array}$, may be inconsistent}
\opt{D}{For some given (a, b, c) $\in$ S, the system of equation $\begin{array}{l}\mathrm{ax+by}=1\\\mathrm{bx+cy}=-1\end{array}$, has a unique solution}
\end{qoptions}
\qn{7} Let $\mathbb R^3$ denote the three-dimensional space. Take two points P = (3, 4, 5) \& Q(6, 4, 9). Let dist (X, Y) denote distance between two point X and Y in R$^3$.\par
Let S = \{X $\in\mathbb R^3$ : (dist (X, P))$^2$ $-$ (dist (X, Q))$^2$ = 100\} and T = \{Y $\in\mathbb R^3$ : (dist (Y, Q))$^2$ $-$ (dist (Y, P))$^2$ = 100\}, then which of the following is True ?
\begin{qoptions}
\opt{A}{There is a triangle whose area is 2 sq. unit \& all vertices are from S.}
\opt{B}{There are two distinct points L and M in T such that each point on the line segment LM is also in T.}
\opt{C}{There are infinitely many rectangles of perimeter 48, two of whose vertices are from S and the other two vertices are from T.}
\opt{D}{There is a square of perimeter 48, two of whose vertices are from S and the other two vertices are from T.}
\end{qoptions}
\sect{SECTION-I(iii)}
\qn{8} Let $\alpha$, $\beta$ and $\gamma$ be the distinct roots of the equation x$^3$ $-$ 7x + 6 = 0. Consider the set T = \{$\alpha$, $\beta$, $\gamma$\} for a 3 $\times$ 3 matrix M = (a$_{\mathrm{ij}}$)$_{3\times3}$, define R$_\mathrm i$ = a$_{\mathrm{i1}}$ + a$_{\mathrm{i2}}$ + a$_{\mathrm{i3}}$ and C$_\mathrm j$ = a$_{\mathrm{1j}}$ + a$_{\mathrm{2j}}$ + a$_{\mathrm{3j}}$ for i = 1, 2, 3 and j = 1, 2, 3.\par
Match each entry in \textbf{List-I} to the correct entry in \textbf{List-II}.\par
{\footnotesize\setlength{\tabcolsep}{1.5pt}\noindent\begin{tabular}{|c|L{5.4cm}|c|L{1.3cm}|}
\hline
 & \multicolumn{1}{>{\centering\arraybackslash}p{5.4cm}|}{List-I} & & \multicolumn{1}{>{\centering\arraybackslash}p{1.3cm}|}{List-II} \\ \hline
(P) & The number of matrices M = (a$_{\mathrm{ij}}$)$_{3\times3}$ with all entries in T such that R$_\mathrm i$ = C$_\mathrm j$ = 0 for all i, j, is & (1) & 1 \\ \hline
(Q) & The number of symmetric matrices M = (a$_{\mathrm{ij}}$)$_{3\times3}$ with all entries in T such that C$_\mathrm j$ = 0 for all j, is & (2) & 12 \\ \hline
(R) & Let M = (a$_{\mathrm{ij}}$)$_{3\times3}$ be a skew symmetric matrix such that a$_{\mathrm{ij}}\in$ T for i $>$ j. Then the number of elements in the set $\bigl\{\begin{pmatrix}x\\y\\z\end{pmatrix}:x,y,z\in\mathbb R,$ $\ M\begin{pmatrix}x\\y\\z\end{pmatrix}=\begin{pmatrix}a_{12}\\0\\-a_{23}\end{pmatrix}\bigr\}$ is & (3) & infinite \\ \hline
(S) & Let M = (a$_{\mathrm{ij}}$)$_{3\times3}$ be a matrix with all entries in T such that R$_\mathrm i$ = 0 for all i. Then the absolute value of the determinant of M is & (4) & 6 \\ \hline
 & & (5) & 0 \\ \hline
\end{tabular}}\par
\begin{qoptions}
\opt{A}{P $\to$ 2;Q $\to$ 4;R $\to$ 3;S $\to$ 5}
\opt{B}{P $\to$ 3;Q $\to$ 2;R $\to$ 4;S $\to$ 2}
\opt{C}{P $\to$ 2;Q $\to$ 2;R $\to$ 5;S $\to$ 2}
\opt{D}{P $\to$ 3;Q $\to$ 4;R $\to$ 3;S $\to$ 1}
\end{qoptions}
\qn{9} Let $\alpha\in$ R be such that lines :\par
L$_1$ : $\dfrac{x-1}2=\dfrac{y-1}3=\dfrac{z-4}\alpha$ and\par
L$_2$ : $\dfrac x3=\dfrac y4=\dfrac{z-2}4$\par
intersect let R$_1$ be the point of intersection of L$_1$ \& L$_2$. Let O = (0, 0, 0) and $\hat{\mathrm n}$ be a unit vector normal to plane containing lines L$_1$ \& L$_2$ :\par
{\footnotesize\setlength{\tabcolsep}{1.5pt}\noindent\begin{tabular}{|c|L{2.6cm}|c|L{3.4cm}|}
\hline
 & \multicolumn{1}{>{\centering\arraybackslash}p{2.6cm}|}{List-I} & & \multicolumn{1}{>{\centering\arraybackslash}p{3.4cm}|}{List-II} \\ \hline
(P) & $\alpha$ = & (1) & $\dfrac2{\sqrt{21}}$ \\ \hline
(Q) & A possible choice for $\hat{\mathrm n}$ = & (2) & 4 \\ \hline
(R) & $\overrightarrow{\mathrm{OR}_1}$ = & (3) & $\left(\dfrac{4\hat{\mathrm i}}{\sqrt{21}}-\dfrac{2\hat{\mathrm j}}{\sqrt{21}}-\dfrac{\hat{\mathrm k}}{\sqrt{21}}\right)$ \\ \hline
(S) & $\left|\overrightarrow{\mathrm{OR}_1}\cdot\hat{\mathrm n}\right|$ & (4) & 2 \\ \hline
 & & (5) & 3$\hat{\mathrm i}$ + 4$\hat{\mathrm j}$ + 6$\hat{\mathrm k}$ \\ \hline
\end{tabular}}\par
\begin{qoptions}
\opt{A}{P $\to$ 2;Q $\to$ 3;R $\to$ 4;S $\to$ 2}
\opt{B}{P $\to$ 4;Q $\to$ 1;R $\to$ 5;S $\to$ 3}
\opt{C}{P $\to$ 2;Q $\to$ 1;R $\to$ 4;S $\to$ 1}
\opt{D}{P $\to$ 4;Q $\to$ 3;R $\to$ 5;S $\to$ 1}
\end{qoptions}
\qn{10} Let line 2y = x + 3 touch a circle with centre ($\alpha$, 0), ($\alpha$ $>$ 0) and radius r at a point A$_1$. Let B$_1$ be the point on the circle such that the line segment A$_1$B$_1$ is a diameter of the circle. Let $\alpha$ + r = 2 + $\sqrt5$. Match List-I to List-II.\par
{\footnotesize\setlength{\tabcolsep}{1.5pt}\noindent\begin{tabular}{|c|c|c|c|}
\hline
\multicolumn{2}{|c|}{List-I} & \multicolumn{2}{c|}{List-II} \\ \hline
(P) & $\alpha$ = & (1) & (1, 2) \\ \hline
(Q) & r = & (2) & $\sqrt5$ \\ \hline
(R) & A$_1$ = & (3) & 2 \\ \hline
(S) & B$_1$ = & (4) & (3, $-$2) \\ \hline
 & & (5) & 1 \\ \hline
\end{tabular}}\par
\begin{qoptions}
\opt{A}{P $\to$ 3;Q $\to$ 5;R $\to$ 4;S $\to$ 1}
\opt{B}{P $\to$ 5;Q $\to$ 3;R $\to$ 1;S $\to$ 2}
\opt{C}{P $\to$ 3;Q $\to$ 2;R $\to$ 1;S $\to$ 4}
\opt{D}{P $\to$ 5;Q $\to$ 2;R $\to$ 3;S $\to$ 4}
\end{qoptions}
\qn{11} Let f : R $\to$ R \& g : R $\to$ R, be functions defined by f(x) = $\begin{cases}x^2\sin\frac1x,&x\neq0\\0,&x=0\end{cases}$ \& g(x) = $\begin{cases}1-x,&0\leqslant x<1\\0,&\text{otherwise}\end{cases}$ where a, b, c, d $\in$ R. Let h(x) be a function, h : R $\to$ R, h(x) = af(x) + b(g(x) + g(1 $-$ x)) + c(x $-$ g(x)) + dg(x), x $\in$ R\par
Match List-I \& List-II\par
{\footnotesize\setlength{\tabcolsep}{1.5pt}\noindent\begin{tabular}{|c|L{3.6cm}|c|L{3.0cm}|}
\hline
 & \multicolumn{1}{>{\centering\arraybackslash}p{3.6cm}|}{List-I} & & \multicolumn{1}{>{\centering\arraybackslash}p{3.0cm}|}{List-II} \\ \hline
(P) & If a = 0, b = 1, c = 0 \& d = 0 & (1) & h is one-one \\ \hline
(Q) & If a = 1, b = 0, c = 0, d = 0 & (2) & h is onto \\ \hline
(R) & If a = 0, b = 0, c = 1, d = 0 & (3) & h is differentiable \\ \hline
(S) & If a = 0, b = 0, c = 0 \& d = 1 & (4) & Range of h is [0, 1] \\ \hline
 & & (5) & Range of h is \{0, 1\} \\ \hline
\end{tabular}}\par
\begin{qoptions}
\opt{A}{P $\to$ 4;Q $\to$ 1;R $\to$ 3;S $\to$ 4}
\opt{B}{P $\to$ 4;Q $\to$ 2;R $\to$ 1;S $\to$ 5}
\opt{C}{P $\to$ 5;Q $\to$ 3;R $\to$ 2;S $\to$ 4}
\opt{D}{P $\to$ 5;Q $\to$ 2;R $\to$ 1;S $\to$ 5}
\end{qoptions}
\sect{SECTION-II}
\qn{1} Let a = 2$\sqrt3$ \& b = $\sqrt[3]{3\sqrt2}$\par
2x + 3y = log$_\mathrm a$144 + log$_\mathrm b\sqrt{18}$\par
2x $-$ y = 2log$_6$ab$^3$\par
then value of 3x + 4y is equal to \underline{\hspace{1.5cm}}
\qn{2} Let f(x) = x$^4$ + ax$^3$ + bx$^2$ + c be a polynomial with real coefficients such that f(1) = $-$7 suppose that $\sqrt2$i is a root of the equation 4x$^3$ + 3ax$^2$ + 2bx = 0, where i = $\sqrt{-1}$. If $\alpha_1$, $\alpha_2$, $\alpha_3$ and $\alpha_4$ are all the roots of the equation f(x) = 0, then $|\alpha_1|^2+|\alpha_2|^2+|\alpha_3|^2+|\alpha_4|^2$ is equal to \underline{\hspace{1.5cm}}
\qn{3} A group of 12 students S$_1$, S$_2$, S$_3$ ..., S$_{12}$ are to be divided into three teams X, Y, Z of sizes 3, 4 \& 5 respectively. If S$_1$ \& S$_2$ can not go to team X \& S$_3$ can not go to team Y, then number of ways to form such teams is `N', then $\dfrac{\mathrm N}2$ is equal to \underline{\hspace{1.5cm}}
\qn{4} Let $\overrightarrow{\mathrm{OP}}$ = ($\alpha-1$)$\hat{\mathrm i}$ + $\hat{\mathrm j}$ + 2$\hat{\mathrm k}$, $\overrightarrow{\mathrm{OQ}}$ = ($\beta-2$)$\hat{\mathrm i}$ + 2$\hat{\mathrm j}$ $-$ $\hat{\mathrm k}$ and $\overrightarrow{\mathrm{OR}}$ = $\hat{\mathrm i}$ + $\hat{\mathrm j}$ + $\hat{\mathrm k}$ be three vectors, where a, b $\in$ R $-$ \{0\} and O denotes the origin. If $\left(\overrightarrow{\mathrm{OP}}\times\overrightarrow{\mathrm{OQ}}\right)\cdot\overrightarrow{\mathrm{OR}}$ = 0 and the point ($\alpha$, $\beta$, 2) lies on the plane 3x + y $-$ z $-$ $\ell$ = 0, then the value of $\ell$ is \underline{\hspace{1.5cm}}
\qn{5} Let S = $\bigl\{\mathrm A=\begin{pmatrix}0&1&\mathrm c\\1&2\mathrm a&\mathrm d\\1&2\mathrm b&\mathrm e\end{pmatrix}:\mathrm a,\mathrm b,\mathrm c,\mathrm d,\mathrm e\in\{0,1\}$ and $|\mathrm A|\in\{-2,2\}\bigr\}$, where $|$A$|$ denotes the determinant of A.\par
Then the number of elements is S is \underline{\hspace{1.5cm}}
\qn{6} Let X be a random variable and let P(X = x) denote the probability that X takes the value x. Suppose that the points (x, P(X = x)), x = 1, 2, 3, 4, 5 lie on a fixed straight line in the xy-plane and P(X = x) = 0 for all x $\in$ R $-$ \{1, 2, 3, 4, 5\}. If the mean of X is $\dfrac{10}3$ and the variance of X is $\alpha$, then the value of 9$\alpha$ is \underline{\hspace{1.5cm}}
\raggedcolumns
\papersec{2}
\partsec{1}{PHYSICS}
\sect{SECTION-I(i)}
\qn{1} A region in the form of an equilateral triangle (in x-y plane) of height \textit L has a uniform magnetic field $\vec{\mathrm B}$ pointing in the +z-direction. A conducting loop PQR, in the form of an equilateral triangle of the same height \textit L, is placed in the \textit{x-y} plane with its vertex P at \textit x = 0 in the orientation shown in the figure. At \textit t = 0, the loop starts entering the region of the magnetic field with a uniform velocity $\vec{\mathrm v}$ along the +\textit x-direction. The plane of the loop and its orientation remain unchanged throughout its motion.\par
\qfig{p2_phy_s1_q01_1}
Which of the following graph best depicts the variation of the induced emf (\textit E) in the loop as a function of the distance (\textit x) starting from \textit x = 0 ?
\begin{qoptions}
\opt{A}{\qfigi{p2_phy_s1_q01_A}}
\opt{B}{\qfigi{p2_phy_s1_q01_B}}
\opt{C}{\qfigi{p2_phy_s1_q01_C}}
\opt{D}{\qfigi{p2_phy_s1_q01_D}}
\end{qoptions}
\qn{2} A particle of mass \textit m is under the influence of the gravitational field of a body of mass \textit M ($\gg$ \textit m). The particle is moving in a circular orbit of radius \textit r$_0$ with time period \textit T$_0$ around the mass \textit M. Then, the particle is subjected to an additional central force, corresponding to the potential energy \textit V$_c$(\textit r) = \textit m$\alpha$/\textit r$^3$, where $\alpha$ is a positive constant of suitable dimensions and \textit r is the distance from the center of the orbit. If the particle moves in the same circular orbit of radius \textit r$_0$ in the combined gravitational potential due to \textit M and \textit V$_c$, but with a new time period \textit T$_1$, then $\dfrac{\mathrm T_1^2-\mathrm T_0^2}{\mathrm T_1^2}$ is given by :- [\textit G is the gravitational constant]
\begin{qoptions}
\opt{A}{$\dfrac{\alpha}{2\mathrm{GMr}_0^2}$}
\opt{B}{$\dfrac{2\alpha}{\mathrm{GMr}_0^2}$}
\opt{C}{$\dfrac{\alpha}{\mathrm{GMr}_0^2}$}
\opt{D}{$\dfrac{3\alpha}{\mathrm{GMr}_0^2}$}
\end{qoptions}
\qn{3} A metal target with atomic number \textit Z = 91 is bombarded with a high energy electron beam. The emission of X-rays from the target is analyzed. The ratio r of the wavelengths of the \textit K$_\alpha$-line and the cut-off is found to be \textit r = 8. If the same electron beam bombards another metal target with \textit Z = 41, the value of \textit r will be
\begin{qoptions}
\opt{A}{12.64}
\opt{B}{40.5}
\opt{C}{0.63}
\opt{D}{1.58}
\end{qoptions}
\qn{4} A thin stiff insulated metal wire is bent into a circular loop with its two ends extending tangentially from the same point of the loop. The wire loop has mass \textit m and radius \textit r and it is in a uniform vertical magnetic field \textit B$_0$, as shown in the figure. Initially, it hangs vertically downwards, because of acceleration due to gravity \textit g, on two conducting supports at P and Q. When a current \textit I is passed through the loop, the loop turns about the line PQ by maximum angle $\theta$ given by :-\par
\qfig{p2_phy_s1_q04_1}
\begin{qoptions}
\opt{A}{tan $\theta$ = $\dfrac{\pi\mathrm{rIB}_0}{\mathrm{mg}}$}
\opt{B}{tan $\theta$ = $\dfrac{2\pi\mathrm{rIB}_0}{\mathrm{mg}}$}
\opt{C}{tan $\dfrac\theta2$ = $\dfrac{\pi\mathrm{rIB}_0}{\mathrm{mg}}$}
\opt{D}{tan $\dfrac\theta2$ = $\dfrac{2\pi\mathrm{rIB}_0}{\mathrm{mg}}$}
\end{qoptions}
\sect{SECTION-I(ii)}
\qn{5} A small electric dipole $\vec{\mathrm p}_0$, having a moment of inertia I about its center, is kept at a distance \textit r from the center of a spherical shell of radius \textit R. The surface charge density $\sigma$ is uniformly distributed on the spherical shell. The dipole is initially oriented at a small angle $\theta$ as shown in the figure. While staying at a distance \textit r, the dipole is free to rotate about its center. If released from rest, then which of the following statement(s) is(are) correct? [$\epsilon_0$ is the permittivity of free space.]\par
\qfig{p2_phy_s2_q05_1}
\begin{qoptions}
\opt{A}{The dipole will undergo small oscillations at any finite value of \textit r.}
\opt{B}{The dipole will undergo small oscillations at any finite value of \textit r $>$ \textit R.}
\opt{C}{The dipole will undergo small oscillations with an angular frequency of $\sqrt{\dfrac{\sigma\mathrm p_0}{4\epsilon_0\mathrm I}}$ at \textit r = 2\textit R.}
\opt{D}{The dipole will undergo small oscillations with an angular frequency of $\sqrt{\dfrac{\sigma\mathrm p_0}{10\epsilon_0\mathrm I}}$ at \textit r = 10\textit R.}
\end{qoptions}
\qn{6} A table tennis ball has radius (3/2) $\times$ 10$^{-2}$ m and mass (22/7) $\times$ 10$^{-3}$ kg. It is slowly pushed down into a swimming pool to a depth of \textit d = 0.7 m below the water surface and then released from rest. It emerges from the water surface at speed \textit v, without getting wet, and rises up to a height \textit H. Which of the following option(s) is(are) correct? [Given: $\pi$ = 22/7, \textit g = 10 ms$^{-2}$, density of water = 1 $\times$ 10$^3$ kg m$^{-3}$, viscosity of water = 1 $\times$ 10$^{-3}$ Pa-s.]
\begin{qoptions}
\opt{A}{The work done in pushing the ball to the depth \textit d is 0.077 J.}
\opt{B}{If we neglect the viscous force in water, then the speed \textit v = 7 m/s.}
\opt{C}{If we neglect the viscous force in water, then the height \textit H = 1.4 m.}
\opt{D}{The ratio of the magnitudes of the net force excluding the viscous force to the maximum viscous force in water is 250.}
\end{qoptions}
\qn{7} A positive, singly ionized atom of mass number \textit A$_\mathrm M$ is accelerated from rest by the voltage 192 V. Thereafter, it enters a rectangular region of width \textit w with magnetic field $\vec{\mathrm B}_0$ = 0.1$\hat{\mathrm k}$ Tesla, as shown in the figure. The ion finally hits a detector at the distance x below its starting trajectory. Which of the following option(s) is(are) correct? [Given: Mass of neutron/proton = (5/3) $\times$ 10$^{-27}$ kg, charge of the electron = 1.6 $\times$ 10$^{-19}$ C.]\par
\qfig{p2_phy_s2_q07_1}
\begin{qoptions}
\opt{A}{The value of x for \textit H$^+$ ion is 4 cm.}
\opt{B}{The value of x for an ion with \textit A$_\mathrm M$ = 144 is 48 cm.}
\opt{C}{For detecting ions with 1 $\leq$ \textit A$_\mathrm M$ $\leq$ 196, the minimum height (x$_1$ $-$ x$_0$) of the detector is 52 cm.}
\opt{D}{The minimum width \textit w of the region of the magnetic field for detecting ions with \textit A$_\mathrm M$ = 196 is 28 cm.}
\end{qoptions}
\sect{SECTION-II(i)}
\paraq{1}{2}
In a Young's double slit experiment, each of the two slits A and B, as shown in the figure, are oscillating about their fixed center and with a mean separation of 0.8 mm. The distance between the slits at time \textit t is given by \textit d = (0.8 + 0.04 sin $\omega t$) mm, where $\omega$ = 0.08 rad s$^{-1}$. The distance of the screen from the slits is 1 m and the wavelength of the light used to illuminate the slits is 5000 \AA. The interference pattern on the screen changes with time, while the central bright fringe (zeroth fringe) remains fixed at point O.\par
\qfig{p2_phy_s3_para1_2_1}
\qn{1} The 8$^{\mathrm{th}}$ bright fringe above the point O oscillates with time between two extreme positions. The separation between these two extreme positions, in micrometer ($\mu$m), is \underline{\hspace{1.5cm}}.
\qn{2} The maximum speed in $\mu$ m/s at which the 8$^{\mathrm{th}}$ bright fringe will move is \underline{\hspace{1.5cm}}.
\paraq{3}{4}
Two particles, 1 and 2, each of mass \textit m, are connected by a massless spring, and are on a horizontal frictionless plane, as shown in the figure. Initially, the two particles, with their center of mass at \textit x$_0$, are oscillating with amplitude \textit a and angular frequency $\omega$. Thus, their positions at time \textit t are given by\par
\textit x$_1$(\textit t) = (\textit x$_0$ + \textit d) + \textit a sin $\omega t$ and \textit x$_2$(\textit t) = (\textit x$_0$ $-$ \textit d) $-$ \textit a sin $\omega t$, respectively, where \textit d $>$ 2\textit a. Particle 3 of mass \textit m moves towards this system with speed \textit u$_0$ = \textit a$\omega$/2, and undergoes instantaneous elastic collision with particle 2, at time \textit t$_0$. Finally, particles 1 and 2 acquire a center of mass speed v$_\mathrm{cm}$ and oscillate with amplitude \textit b and the same angular frequency $\omega$.\par
\qfig{p2_phy_s3_para3_4_1}
\qn{3} If the collision occurs at time \textit t$_0$ = $\pi$/2$\omega$, the value of \textit v$_\mathrm{cm}$/(\textit a$\omega$) will be \underline{\hspace{1.5cm}}.
\qn{4} If the collision occurs at time \textit t$_0$ = 0, then the value of 4\textit b$^2$/a$^2$ will be \underline{\hspace{1.5cm}}.
\sect{SECTION-II(ii)}
\qn{5} The dimensions of a cone are measured using a scale with a least count of 1 mm. The diameter of the base and the height are both measured to be 10.0 cm. The maximum percentage error in the determination of the volume is \underline{\hspace{1.5cm}}.
\qn{6} A ball is thrown from the location (x$_0$, y$_0$) = (0,0) of a horizontal playground with an initial speed \textit v$_0$ at an angle $\theta_0$ from the +x-direction. The ball is to be hit by a stone, which is thrown at the same time from the location (x$_1$, y$_1$) = (\textit L, 0). The stone is thrown at an angle (180 $-$ $\theta_1$) from the +x-direction with a suitable initial speed. For a fixed \textit v$_0$, when ($\theta_0$, $\theta_1$) = (45$^\circ$, 45$^\circ$), the stone hits the ball after time \textit T$_1$, and when ($\theta_0$, $\theta_1$) = (60$^\circ$, 30$^\circ$), it hits the ball after time \textit T$_2$. In such a case, (\textit T$_1$/\textit T$_2$)$^4$ is \underline{\hspace{1.5cm}}.
\qn{7} A charge is kept at the central point P of a cylindrical region. The two edges subtend a half-angle $\theta$ at P, as shown in the figure. When $\theta$ = 60$^\circ$, then the electric flux through the curved surface of the cylinder is $\Phi$. If $\theta$ = 30$^\circ$, then the electric flux through the curved surface becomes $\sqrt{\mathrm n}\phi$, where the value of n is \underline{\hspace{1.5cm}}.\par
\qfig{p2_phy_s4_q07_1}
\qn{8} Two equilateral-triangular prisms P$_1$ and P$_2$ are kept with their sides parallel to each other and very close in vacuum, as shown in the figure. A light ray enters prism P$_1$ at an angle of incidence $\theta$ such that the outgoing ray undergoes maximum deviation in prism P$_2$. If the respective refractive indices of P$_1$ and P$_2$ are $\dfrac53$ and 2, then $\theta$ = sin$^{-1}\left(\dfrac53\sin(\beta)\right)$, where the value of $\beta$ (in degree) is \underline{\hspace{1.5cm}}.\par
\qfig{p2_phy_s4_q08_1}
\qn{9} An infinitely long thin wire, having a uniform charge density per unit length of 5 nC/m, is passing through a spherical shell of radius 1 m, as shown in the figure. A 2 nC charge is distributed uniformly over the spherical shell. If the configuration of the charges remains static, the magnitude of the potential difference between points P and R, in Volt, is \underline{\hspace{1.5cm}}. [Given: In SI units $\dfrac1{4\pi\epsilon_0}$ = 9 $\times$ 10$^9$, ln 2 = 0.7. Ignore the area pierced by the wire.]\par
\qfig{p2_phy_s4_q09_1}
\qn{10} A spherical soap bubble inside an air chamber at pressure \textit P$_0$ = 10$^5$ Pa has a certain radius so that the excess pressure inside the bubble is $\Delta P$ = 36 Pa. Now, the chamber pressure is reduced to 8\textit P$_0$/27 so that the bubble radius and its excess pressure change. In this process, all the temperatures remain unchanged. Assume air to be an ideal gas and the excess pressure $\Delta P$ in both the cases to be much smaller than the chamber pressure. The new excess pressure $\Delta P$ in Pa is \underline{\hspace{1.5cm}}.
\raggedcolumns
\papersec{2}
\partsec{2}{CHEMISTRY}
\sect{SECTION-I(i)}
\qn{1} According to Bohr's model, the lowest kinetic energy is associated with the electron in the :
\begin{qoptions}
\opt{A}{first orbit of H atom}
\opt{B}{first orbit of He$^+$}
\opt{C}{third orbit of He$^+$}
\opt{D}{second orbit of Li$^{2+}$}
\end{qoptions}
\qn{2} In a metal deficient oxide sample, \textbf M$_\mathbf x$\textbf Y$_2$O$_4$ (\textbf M and \textbf Y are metals), \textbf M is present in both +2 and +3 oxidation states and \textbf Y is in +3 oxidation state. If the fraction of \textbf M$^{\mathbf{2+}}$ ions present in \textbf M is $\dfrac34$, the value of \textbf X is \underline{\hspace{1.2cm}}.
\begin{qoptions}
\opt{A}{0.25}
\opt{B}{0.89}
\opt{C}{0.67}
\opt{D}{0.75}
\end{qoptions}
\qn{3} In the following reaction sequences, the major product \textbf R is:\par
\qfig{p2_chem_s1_q03_1}
\begin{qoptions}
\opt{A}{\qfigi{p2_chem_s1_q03_A}}
\opt{B}{\qfigi{p2_chem_s1_q03_B}}
\opt{C}{\qfigi{p2_chem_s1_q03_C}}
\opt{D}{\qfigi{p2_chem_s1_q03_D}}
\end{qoptions}
\qn{4} At room temperature, disproportionation of an aqueous solution of \textit{in situ} generated nitrous acid (HNO$_2$) gives the species
\begin{qoptions}
\opt{A}{H$_3$O$^+$, NO$_3^-$ and N$_2$O}
\opt{B}{H$_3$O$^+$, NO$_3^-$ and NO$_2$}
\opt{C}{H$_3$O$^+$, NO$^-$ and NO$_2$}
\opt{D}{H$_3$O$^+$, NO$_3^-$ and NO}
\end{qoptions}
\sect{SECTION-I(ii)}
\qn{5} An aqueous solution of hydrazine (N$_2$H$_4$) is electrochemically oxidized by O$_2$, thereby releasing chemical energy in the form of electrical energy. One of the products generated from the electrochemical reaction is N$_2$(g). Choose the correct statement(s) about the above process :
\begin{qoptions}
\opt{A}{OH$^-$ ions react with N$_2$H$_4$ at the anode to form N$_2$(g) and water, releasing 4 electrons to the anode.}
\opt{B}{At the cathode, N$_2$H$_4$ breaks to N$_2$(g) and nascent hydrogen released at the electrode reacts with oxygen to form water.}
\opt{C}{At the cathode, molecular oxygen gets converted to OH$^-$.}
\opt{D}{Oxides of nitrogen are major by-products of the electrochemical process.}
\end{qoptions}
\qn{6} The option(s) with correct sequence of reagents for the conversion of \textbf P to \textbf Q is(are) :\par
\qfig{p2_chem_s2_q06_1}
\begin{qoptions}
\opt{A}{i) NaBH$_4$ ; ii) SnCl$_2$/HCl ; iii) H$_3$O$^+$ ; iv) Lindlar's catalyst, H$_2$}
\opt{B}{i) Lindlar's catalyst, H$_2$ ; ii) H$_3$O$^+$ ; iii) SnCl$_2$/HCl ; iv) NaBH$_4$}
\opt{C}{i) Lindlar's catalyst, H$_2$ ; ii) SnCl$_2$/HCl ; iii) NaBH$_4$ ; iv) H$_3$O$^+$}
\opt{D}{i) Lindlar's catalyst, H$_2$ ; ii) NaBH$_4$ ; iii) SnCl$_2$/HCl ; iv) H$_3$O$^+$}
\end{qoptions}
\qn{7} Among the following options, select the option(s) in which each complex in \textbf{Set-I} shows geometrical isomerism and the two complexes in \textbf{Set-II} are ionization isomers of each other.\par
en = H$_2$NCH$_2$CH$_2$NH$_2$ ; \qfigi{p2_chem_s2_q07_1}
\begin{qoptions}
\opt{A}{\textbf{Set-I} : [Pt(NH$_3$)$_2$Cl$_2$] and [PdCl$_2$(PPh$_3$)$_2$]\par \textbf{Set-II} : [Co(en)$_2$(NH$_3$)NO$_3$] SO$_4$ and [Co(en)$_2$(NH$_3$SO$_4$] NO$_3$}
\opt{B}{\textbf{Set-I} : [Co(en)(NH$_3$)$_2$Cl$_2$] and [PdCl$_2$(PPh$_3$)$_2$]\par \textbf{Set-II} : [Co(NH$_3$)$_6$] [Cr(CN)$_6$] and [Cr(NH$_3$)$_6$] [Co(CN)$_6$]}
\opt{C}{\textbf{Set-I} : [Pt(py)$_2$\allowbreak(NH$_3$)$_2$\allowbreak Cl$_2$]\allowbreak Br$_2$ and [Co(en)Br$_2$\allowbreak(NH$_3$)$_2$]\par \textbf{Set-II} : [Co(NH$_3$)$_5$\allowbreak Cl]\allowbreak SO$_4$ and [Co(NH$_3$)$_5$\allowbreak(SO$_4$)]\allowbreak Cl}
\opt{D}{\textbf{Set-I} : [Cr(NH$_3$)$_4$Cl$_2$]Cl and [Co(pn)$_3$]Cl$_3$\par \textbf{Set-II} : [Pt(NH$_3$)$_4$Cl$_2$]Br$_2$ and [Pt(NH$_3$)$_4$Br$_2$]Cl$_2$}
\end{qoptions}
\sect{SECTION-II(i)}
\paraq{1}{2}
An organic compound \textbf P with molecular formula C$_9$H$_{18}$O$_2$ decolorizes bromine water and also shows positive iodoform test. \textbf P on ozonolysis followed by treatment with H$_2$O$_2$ gives \textbf Q and \textbf R. While compound \textbf Q shows positive iodoform test, compound \textbf R does not give positive iodoform test. \textbf Q and \textbf R on oxidation with pyridinium chlorochromate (PCC) followed by heating give \textbf S and \textbf T, respectively. Both \textbf S and \textbf T show positive iodoform test.\par
Complete copolymerization of 510 moles of \textbf Q and 510 moles of \textbf R gives one mole of a single acyclic copolymer \textbf U.\par
[Given, atomic mass : H = 1, C =12, O = 16]
\qn{1} Sum of number of sp$^3$ carbon atoms in \textbf S and \textbf T is \underline{\hspace{1.5cm}}.
\qn{2} The value of $\dfrac{\mathrm U}{18}$ is \underline{\hspace{1.5cm}}\par
where \textbf U is molecular weight of copolymer.
\paraq{3}{4}
When potassium iodide is added to an aqueous solution of potassium ferricyanide, a reversible reaction is observed in which a complex \textbf P is formed. In a strong acidic medium, the equilibrium shifts completely towards \textbf P. Addition of ferric chloride to \textbf P in a slightly acidic medium results in insoluble complex \textbf Q.
\qn{3} If 3 moles of potassium iodide are taken then X moles of P and Y moles of iodine is formed. Find the value of (X+4Y).
\qn{4} The number of iron present in the molecular formula of \textbf Q is \underline{\hspace{1.5cm}}.
\sect{SECTION-II(ii)}
\qn{5} To form a complete monolayer of acetic acid on 1g of charcoal, 100 mL of 0.6 M acetic acid was used. Some of the acetic acid ramained unadsorbed. To neutralize the unadsorbed acetic acid, 40 mL of 1 M NaOH solution was required. If each molecule of acetic acid occupies \textbf P $\times$ 10$^{-23}$ m$^2$ surface area on charcoal, the value of \textbf P is \underline{\hspace{1.5cm}}.\par
[Use given data : Surface area of charcoal = 1.5 $\times$ 10$^2$ m$^2$g$^{-1}$; Avogadro's number (N$_\mathrm A$) = 6.0 $\times$ 10$^{23}$ mol$^{-1}$]
\qn{6} Vessel-1 contains \textbf w$_2$ g of a non-volatile solute \textbf X dissolved in \textbf w$_1$ g of water. Vessel-2 contains \textbf w$_2$ g of another non-volatile solute \textbf Y dissolved in \textbf w$_1$ g of water. Both the vessels are at the same temperature and pressure. The molar mass of \textbf X is 60\% of that of \textbf Y. The van't Hoff factor for \textbf X is 20\% more than that of \textbf Y for their respective concentrations.\par
The elevation of boiling point for solution in Vessel-1 is \underline{\hspace{1.5cm}}\% of the solution in Vessel-2.
\qn{7} For a double strand DNA, one strand is given below:\par
\qfig{p2_chem_s4_q07_1}
The amount of energy required to split the double strand DNA into two single strands is \underline{\hspace{1.5cm}}kcal mol$^{-1}$.\par
[Given : Average energy per H-bond for A-T base pair = 1.0 kcal mol$^{-1}$, G-C base pair = 1.5 kcal mol$^{-1}$, and A-U base pair = 1.25 kcal mol$^{-1}$. Ignore electrostatic repulsion between the phosphate groups.]
\qn{8} A sample initially contains only U-238 isotope of uranium. With time, some of the U-238 radioactively decays into Pb-206 while the rest of it remains undisintegrated.\par
When the age of the sample is \textbf P $\times$ 10$^8$ years, the ratio of moles of Pb-206 to that of U-238 in the sample is found to be 15. The value of \textbf P is \underline{\hspace{1.5cm}}.\par
[Given : Half-life of U-238 is 4.5 $\times$ 10$^9$ years]
\qn{9} Among [Co(CN)$_4$]$^{4-}$, [Co(CO)$_3$(NO)], XeF$_4$, [PCl$_4$]$^+$, SF$_4$ , [ICl$_4$]$^-$, [Cu(CN)$_4$]$^{3-}$ , [Ni(CO)$_4$] , [CdCl$_4$]$^{2-}$ , C$_2$H$_4$ and P$_4$ the total number of species with tetrahedral geometry is \underline{\hspace{1.5cm}}.
\qn{10} An organic compound \textbf P having molecular formula C$_6$H$_6$O$_3$ gives ferric chloride test and does not have intramolecular hydrogen bond. The compound \textbf P reacts with 3 equivalents of NH$_2$OH to produce oxime \textbf Q. Treatment of \textbf P with excess methyl iodide in the presence of KOH produces compound \textbf R as the major product. Reaction of \textbf R with excess \textit{methyl}magnesium bromide followed by treatment with H$_3$O$^+$ gives compound \textbf S as the major product.\par
The total number of methyl ($-$CH$_3$) group(s) in compound \textbf S is \underline{\hspace{1.5cm}}.
\raggedcolumns
\papersec{2}
\partsec{3}{MATHEMATICS}
\sect{SECTION-I(i)}
\qn{1} Considering only the principal values of the inverse trigonometric functions, the value of\par
$\tan\left(\sin^{-1}\left(\dfrac45\right)-2\cos^{-1}\left(\dfrac2{\sqrt5}\right)\right)$ is
\begin{qoptions}
\opt{A}{$\dfrac7{24}$}
\opt{B}{$\dfrac{-7}{24}$}
\opt{C}{$\dfrac{-5}{24}$}
\opt{D}{0}
\end{qoptions}
\qn{2} Let S = $\bigl\{(\mathrm x,\mathrm y)\in\mathbb R\times\mathbb R:\mathrm x\geqslant0,$ $\mathrm y\geqslant0,$ $\mathrm y^2\leqslant4\mathrm x,$ $\mathrm y^2\leqslant12-2\mathrm x$ and $3\mathrm y+\sqrt8\,\mathrm x\leqslant5\sqrt8\bigr\}$. If the area of the region S is $17\sqrt2(\alpha)$, then $\alpha$ is equal to
\begin{qoptions}
\opt{A}{$\dfrac{17}2$}
\opt{B}{$\dfrac{17}4$}
\opt{C}{$\dfrac13$}
\opt{D}{$\dfrac{17}3$}
\end{qoptions}
\qn{3} Let k $\in\mathbb R^+$. If $\displaystyle\lim_{x\to0^+}\left(\sin\left(\sin\mathrm k^2\mathrm x\right)+\cos\mathrm x-\mathrm x\right)^{\frac2{\mathrm x}}=\mathrm e^6$, then the value of k is
\begin{qoptions}
\opt{A}{1}
\opt{B}{2}
\opt{C}{3}
\opt{D}{4}
\end{qoptions}
\qn{4} Let f : $\mathbb R\to\mathbb R$ be a function defined by f(x) = $\begin{cases}\mathrm x^2\sin\left(\frac\pi{\mathrm x^2}\right),&\text{if x}\neq0\\0,&\text{if x = 0}\end{cases}$. Then which of the following statements is TRUE ?
\begin{qoptions}
\opt{A}{f(x) = 0 has finitely many solutions in the interval $\left[\dfrac1{10^{10}},\infty\right)$}
\opt{B}{f(x) = 0 has no solutions in the interval $\left[\dfrac1\pi,\infty\right)$}
\opt{C}{The set of solutions of f(x) = 0 in the interval $\left(0,\dfrac1{10^{10}}\right)$ is finite}
\opt{D}{f(x) = 0 has no solutions in the interval $\left(\dfrac1{\pi^2},\dfrac1\pi\right)$}
\end{qoptions}
\sect{SECTION-I(ii)}
\qn{5} Let S be the set of all ($\alpha$, $\beta$) $\in\mathbb R\times\mathbb R$ such that $\displaystyle\lim_{x\to\infty}\dfrac{\sin\left(x^2\right)\left(\log_ex\right)^\alpha\tan\left(\frac1{x^3}\right)}{x^{\alpha\beta}\left(\log_e(1+x)\right)^\beta}=0$. Then which of the following is (are) correct ?
\begin{qoptions}
\opt{A}{($-$1, 3) $\in$ S}
\opt{B}{($-$1, 1) $\in$ S}
\opt{C}{(1, $-$1) $\in$ S}
\opt{D}{(1, $-$2) $\in$ S}
\end{qoptions}
\qn{6} A straight line drawn from the point P(1, 3, 2), parallel to the line $\dfrac{x-2}1=\dfrac{y-4}2=\dfrac{z-6}1$, intersects the plane L$_1$ : x $-$ y + 3z = 6 at the point Q. Another straight line which passes through Q and is perpendicular to the plane L$_1$ intersects the plane L$_2$ : 2x $-$ y + z = $-$4 at the point R. Then which of the following statements is (are) TRUE ?
\begin{qoptions}
\opt{A}{The length of the line segment PQ is $\sqrt6$}
\opt{B}{The perimeter of the triangle PQR is $\sqrt{13}+\sqrt6+\sqrt{11}$}
\opt{C}{The centroid of the triangle PQR is $\left(\dfrac43,\dfrac{14}3,\dfrac53\right)$}
\opt{D}{The coordinates of R are (1, 6, 3)}
\end{qoptions}
\qn{7} Let A$_1$, B$_1$, C$_1$ be three points in the xy-plane. Suppose that the lines A$_1$C$_1$ and B$_1$C$_1$ are tangents to the curve y$^2$ = 8x at A$_1$ and B$_1$, respectively. If O = (0, 0) and C$_1$ = ($-$8, 0), then which of the following statements is (are) TRUE ?
\begin{qoptions}
\opt{A}{The length of the line segment OA$_1$ is 8$\sqrt2$}
\opt{B}{The length of the line segment A$_1$B$_1$ is 16}
\opt{C}{The orthocentre of the triangle A$_1$B$_1$C$_1$ is (0, 0)}
\opt{D}{The orthocentre of the triangle A$_1$B$_1$C$_1$ is ($-$8, 0)}
\end{qoptions}
\sect{SECTION-II(i)}
\paraq{1}{2}
Let S = \{1, 2, 3, 4, 5, 6\} and X be the set of all relations R from S to S that satisfy both the following properties :\par
i. \quad R has exactly 6 elements.\par
ii. \quad For each (a, b) $\in$ R, we have $|$a $-$ b$|\geq$ 2.\par
Let Y = \{R $\in$ X : The range of R has exactly one element\} and\par
Z = \{R $\in$ X : R is a function from S to S\}\par
Let n(A) denote the number of elements in a set A.
\qn{1} If n(X) = $^{\mathrm m}$C$_\mathrm p$, then the maximum value of m + p is \underline{\hspace{1.5cm}}
\qn{2} If the value of n(Z) $-$ n(Y) is k$^2$, then $\sqrt{|\mathrm k|}$ is \underline{\hspace{1.5cm}}
\paraq{3}{4}
Let f : $\left[0,\dfrac\pi2\right]\to$ [0, 1] be the function defined by f(x) = cos$^2$x and let g : $\left[0,\dfrac\pi2\right]\to[0,\infty)$ be the function defined by g(x) = $\sqrt{\dfrac{\pi\mathrm x}2-\mathrm x^2}$.
\qn{3} The value of $2\displaystyle\int_0^{\pi/2}$f(x) g(x) dx $-$ $\displaystyle\int_0^{\pi/2}$g(x) dx = $\alpha$, then $\alpha$ + 2 is \underline{\hspace{1.5cm}}
\qn{4} The value of $\dfrac{64}{\pi^3}\displaystyle\int_0^{\pi/2}$f(x) g(x) dx is \underline{\hspace{1.5cm}}
\sect{SECTION-II(ii)}
\qn{5} Let f : $\mathbb R\to\mathbb R$ be a function such that f(x + y) = f(x) + f(y) for all x, y $\in\mathbb R$ and g : $\mathbb R\to(0,\infty)$ be a function such that g(x + y) = g(x) g(y) for all x, y $\in\mathbb R$. If f$\left(\dfrac{-3}5\right)$ = 12 and g$\left(\dfrac{-1}3\right)$ = 2, then the value of $\left(\mathrm f\left(\dfrac14\right)+\mathrm g(-2)-9\right)$ g(0) is \underline{\hspace{1.5cm}}
\qn{6} A bag contains N balls out of which 3 balls are white, 6 balls are green and the remaining balls are blue. Assume that the balls are identical otherwise. Three balls are drawn randomly one after the other without replacement. For i = 1, 2, 3, let W$_\mathrm i$, G$_\mathrm i$ and B$_\mathrm i$ denote the events that the ball drawn in the i$^{\mathrm{th}}$ draw is a white ball, green ball and blue ball, respectively. If the probability P (W$_1\cap$G$_2\cap$B$_3$) = $\dfrac1{4\mathrm N}$ and the conditional probability P (B$_3|$W$_1\cap$G$_2$) = $\dfrac18$, then N equals \underline{\hspace{1.5cm}}
\qn{7} Let the function f : $\mathbb R\to\mathbb R$ be defined by f(x) = $\dfrac{\sin\mathrm x}{\mathrm e^{\pi\mathrm x}}\cdot\dfrac{\left(\mathrm x^{2025}+2024\mathrm x+2025\right)}{\left(\mathrm x^2-\mathrm x+5\right)}+\dfrac4{\mathrm e^{\pi\mathrm x}}\cdot\dfrac{\left(\mathrm x^{2025}+2024\mathrm x+2025\right)}{\left(\mathrm x^2-\mathrm x+5\right)}$. Then the number of solutions of f(x) = 0 in $\mathbb R$ is \underline{\hspace{1.5cm}}
\qn{8} Let $\vec{\mathrm p}$ = 2$\hat{\mathrm i}$ + $\hat{\mathrm j}$ + 3$\hat{\mathrm k}$ and $\vec{\mathrm q}$ = $\hat{\mathrm i}$ $-$ $\hat{\mathrm j}$ + $\hat{\mathrm k}$. If for some real numbers $\alpha$, $\beta$ and $\gamma$, we have\par
5$\hat{\mathrm i}$ + 2$\hat{\mathrm j}$ + $\hat{\mathrm k}$ = $\alpha\left(2\vec{\mathrm p}-\vec{\mathrm q}\right)$ + $\beta\left(\vec{\mathrm p}+2\vec{\mathrm q}\right)$ + $\gamma\left(\vec{\mathrm p}\times\vec{\mathrm q}\right)$, then the value of 19$\gamma$ is \underline{\hspace{1.5cm}}
\qn{9} A normal with slope $\dfrac1{\sqrt6}$ is drawn from the point (0, $-\alpha$) to the parabola x$^2$ = $-$4ay, where a $>$ 0. Let L be the line passing through (0, $-\alpha$) and parallel to the directrix of the parabola. Suppose that L intersects the parabola at two points A and B. Let r denote the length of the latus rectum and s denote the square of the length of the line segment AB. If r : s = 1 : 32, then the value of 24a is \underline{\hspace{1.5cm}}
\qn{10} Let the function f : [1, $\infty$) $\to\mathbb R$ be defined by\par
\par\noindent\resizebox{\linewidth}{!}{$\mathrm f(t) = \begin{cases}(-1)^{n+1}2,&\text{if t = 2n$-$1, n}\in\mathbb N\\\dfrac{(2n+1-t)}2\mathrm f(2n-1)+\dfrac{(t-(2n-1))}2\mathrm f(2n+1),&\text{if 2n$-$1 $<$ t $<$ 2n + 1, n}\in\mathbb N\end{cases}$}\par
Define g(x) = $\displaystyle\int_1^{x}$f(t) dt, x $\in$ (1, $\infty$). Let $\alpha$ denote the number of solutions of the equation g(x) = 0 in the interval (1, 10] and $\beta$ = $\displaystyle\lim_{x\to1^+}\dfrac{\mathrm g(\mathrm x)}{\mathrm x-1}$. Then the value of $\alpha$ + $\beta$ is equal to \underline{\hspace{1.5cm}}
\end{multicols}
\end{document}
